\noindent \textit{\textcolor{red!60!black}{
1. Gold clause or sensitivity to leverage\\ 
Ideally, to argue that gold-denominated debt is driving the results, one could want to contrast two firms with similar levels of interest-bearing debt but with variation in the fraction of that debt that is gold-denominated. The authors' main measure relates bonds outstanding with gold clauses relative to total liabilities, and controls for book and market leverage, measured by liabilities to assets. What this paper does not do, and what I think they may not be able to do, is to control for interest-bearing debt to total liabilities, or instead relate gold-denominated debt to interest-bearing debt as their main variable. The reason why I think they may not be able to do it is because the vast majority of bonds were indexed, and bank loans were tiny. What this means is that it is hard to know whether the indexing or the level of interest bearing debt are truly driving the results. One can think of reasons why variation in leverage may matter differentially over time independent of indexation. For example, bond markets were especially depressed in 1933 and 1934, but recovered substantially in 1935 and 1936. This is also when the second New Deal reactivated the economy, albeit temporarily. So it would be important to better disentangle gold-indexed bonds versus other bonds, as described above (and the summary statistics for these two should be listed in the paper). If this cannot be done, the paper needs to be more upfront about this important limitation.
}}
\bigskip

\noindent Your impression is correct. In 1932, about 97\% of the corporate bonds in our sample had the gold clause and these bonds represent 99\% of total bond amount outstanding. In our original submission, we had that information in a footnote, but we agree that such an important piece of information should be featured more prominently in the paper. 

To clarify this point, we now mention it on the first page of the introduction: ``When the United States effectively abandoned the gold standard in April 1933, the dollar quickly depreciated against gold, and the existence of these clauses posed a serious threat to most corporate bond issuers: around 97\% of outstanding corporate bonds were subject to a gold clause.'' We hope that explicitly stating the exact number early in the introduction leaves no room for ambiguity. We have also added details on the number of corporate bonds that contained the gold clause versus those that did not in Table 2 of the revised manuscript.

Motivated by your comment, as well as a related suggestion from another referee, we have revised the definition of our exposure measure $d$. In the original submission, the denominator was total liabilities (which, as you noted, includes non-interest-bearing debt). In the revision, the denominator is the sum of corporate bonds, preferred shares, and bank debt---that is, all long-term obligations involving fixed contractual payments:
\[
d_{j,t} = \frac{\text{Gold-denominated debt}_{j,t}}{\text{Corporate bonds}_{j,t} + \text{Preferred shares}_{j,t} + \text{Bank debt}_{j,t}}.
\]
While corporate bonds and preferred shares are not perfect substitutes, they are similar in that they are both claims to future cash flows that are subject to fixed contractual payments. This definition ensures that variation in $d$ more closely reflects differences in the \textit{composition} of fixed-payment claims, rather than differences in its \textit{level}.

In sum, we have made the 97\% figure more prominent in the introduction, provided detailed bond statistics in Table 2, and revised the definition of $d$ to better isolate variation in gold indexation from variation in the level of interest-bearing debt.